\section{Sheaves of algebraic structures}
\label{section-sheaves-algebraic-structures}

\noindent
In Sheaves, Section \ref{sheaves-section-algebraic-structures}
we introduced a type of algebraic structure to be a pair
$(\mathcal{A}, s)$, where $\mathcal{A}$ is a category,
and $s : \mathcal{A} \to \textit{Sets}$ is a functor such
that
\begin{enumerate}
\item $s$ is faithful,
\item $\mathcal{A}$ has limits and $s$ commutes with limits,
\item $\mathcal{A}$ has filtered colimits and $s$ commutes with them, and
\item $s$ reflects isomorphisms.
\end{enumerate}
For such a type of algebraic structure we saw that a presheaf
$\mathcal{F}$ with values in $\mathcal{A}$ on a space $X$ is a sheaf if and
only if the associated presheaf of sets is a sheaf. Moreover,
we worked out the notion of stalk, and given a continuous map
$f : X \to Y$ we defined adjoint functors pushforward and pullback
on sheaves of algebraic structures which agrees with pushforward
and pullback on the underlying sheaves of sets. In addition extending
a sheaf of algebraic structures from a basis to all opens
of a space, works as expected.

\medskip\noindent
Part of this material still works in the setting of sites and sheaves.
Let $(\mathcal{A}, s)$ be a type of algebraic structure.
Let $\mathcal{C}$ be a site. Let us denote
$\textit{PSh}(\mathcal{C}, \mathcal{A})$,
resp.\ $\Sh(\mathcal{C}, \mathcal{A})$ the category
of presheaves, resp.\ sheaves with values in $\mathcal{A}$ on $\mathcal{C}$.
\begin{itemize}
\item[] ($\alpha$) A presheaf with values in $\mathcal{A}$ is
a sheaf if and only if its underlying presheaf of sets is a sheaf.
See the proof of Sheaves, Lemma \ref{sheaves-lemma-sheaves-structure}.
\item[] ($\beta$) Given a presheaf $\mathcal{F}$ with values in
$\mathcal{A}$ the presheaf ${\mathcal{F}}^\# = (\mathcal{F}^+)^+$
is a sheaf. This is true since the colimits in the sheafification process
are filtered, and even colimits over directed sets (see
Section \ref{section-sheafification}, especially the proof of
Lemma \ref{lemma-sheafification-exact})
and since $s$ commutes with filtered colimits.
\item[] ($\gamma$) We get the following commutative diagram
$$
\xymatrix{
\Sh(\mathcal{C}, \mathcal{A}) \ar@<1ex>[r] \ar[d]_s &
\textit{PSh}(\mathcal{C}, \mathcal{A}) \ar@<1ex>[l]^{{}^\#} \ar[d]^s\\
\Sh(\mathcal{C}) \ar@<1ex>[r] &
\textit{PSh}(\mathcal{C}) \ar@<1ex>[l]
}
$$
\item[] ($\delta$) We have $\mathcal{F} = \mathcal{F}^\#$ if and only if
$\mathcal{F}$ is a sheaf of algebraic structures.
\item[] ($\epsilon$) The functor ${}^\#$ is adjoint to the inclusion functor:
$$
\Mor_{\textit{PSh}(\mathcal{C}, \mathcal{A})}(\mathcal{G}, \mathcal{F})
=
\Mor_{\Sh(\mathcal{C}, \mathcal{A})}(\mathcal{G}^\#, \mathcal{F})
$$
The proof is the same as the proof of
Proposition \ref{proposition-sheafification-adjoint}.
\item[] ($\zeta$) The functor
$\mathcal{F} \mapsto \mathcal{F}^\#$ is left exact.
The proof is the same as the proof of Lemma \ref{lemma-sheafification-exact}.
\end{itemize}

\begin{definition}
\label{definition-pushforward-algebraic-structures}
Let $f : \mathcal{D} \to \mathcal{C}$ be a morphism of sites
given by a functor $u : \mathcal{C} \to \mathcal{D}$.
We define the {\it pushforward} functor for presheaves of algebraic structures
by the rule $u^p\mathcal{F}(U) = \mathcal{F}(uU)$,
and for sheaves of algebraic structures by the same rule, namely
$f_*\mathcal{F}(U) = \mathcal{F}(uU)$.
\end{definition}

\noindent
The problem comes with trying the define the pullback.
The reason is that the colimits defining
the functor $u_p$ in Section \ref{section-functoriality-PSh}
may not be filtered. Thus the axioms above are not enough in
general to define the pullback of a (pre)sheaf of algebraic
structures. Nonetheless, in almost all cases the following
lemma is sufficient to define pushforward, and pullback
of (pre)sheaves of algebraic structures.

\begin{lemma}
\label{lemma-push-pull-good-case}
Suppose the functor $u : \mathcal{C} \to \mathcal{D}$ satisfies
the hypotheses of Proposition \ref{proposition-get-morphism},
and hence gives rise to a morphism of sites
$f : \mathcal{D} \to \mathcal{C}$. In this case
the pullback functor $f^{-1}$ (resp.\ $u_p$) and the pushforward
functor $f_*$ (resp. $u^p$) extend to an adjoint pair of functors on
the categories of sheaves (resp.\ presheaves)  of algebraic structures.
Moreover, these functors commute with taking
the underlying sheaf (resp.\ presheaf) of sets.
\end{lemma}

\begin{proof}
We have defined $f_* = u^p$ above.
In the course of the proof of Proposition \ref{proposition-get-morphism}
we saw that all the colimits used to define $u_p$ are
filtered under the assumptions of the proposition.
Hence we conclude from the definition of a type of
algebraic structure that we may define $u_p$ by exactly
the same colimits as a functor on presheaves of algebraic structures.
Adjointness of $u_p$ and $u^p$ is proved in exactly the
same way as the proof of Lemma \ref{lemma-adjoints-u}.
The discussion of sheafification of presheaves of
algebraic structures above then implies that we may define
$f^{-1}(\mathcal{F}) = (u_p\mathcal{F})^\#$.
\end{proof}

\noindent
We briefly discuss a method for dealing with pullback and
pushforward for a general morphism of sites, and more generally
for any morphism of topoi.

\medskip\noindent
Let $\mathcal{C}$ be a site.
In the case $\mathcal{A} = \textit{Ab}$,
we may think of an abelian (pre)sheaf on $\mathcal{C}$
as a quadruple $(\mathcal{F}, +, 0, i)$.
Here the data are
\begin{enumerate}
\item[(D1)] $\mathcal{F}$ is a sheaf of sets,
\item[(D2)] $+ : \mathcal{F} \times \mathcal{F} \to \mathcal{F}$ is
a morphism of sheaves of sets,
\item[(D3)] $0 : * \to \mathcal{F}$ is a morphism from the
singleton sheaf (see Example \ref{example-singleton-sheaf})
to $\mathcal{F}$, and
\item[(D4)] $i : \mathcal{F} \to \mathcal{F}$ is a morphism of sheaves
of sets.
\end{enumerate}
These data have to satisfy the following axioms
\begin{enumerate}
\item[(A1)] $+$ is associative and commutative,
\item[(A2)] $0$ is a unit for $+$, and
\item[(A3)] $+ \circ (1, i) = 0 \circ (\mathcal{F} \to *)$.
\end{enumerate}
Compare Sheaves, Lemma \ref{sheaves-lemma-abelian-presheaves}.
Let $f : \mathcal{D} \to \mathcal{C}$ be a morphism of sites.
Note that since $f^{-1}$ is exact we have
$f^{-1}* = *$ and
$f^{-1}(\mathcal{F} \times \mathcal{F}) =
f^{-1}\mathcal{F} \times f^{-1}\mathcal{F}$.
Thus we can define $f^{-1}\mathcal{F}$ simply as the quadruple
$(f^{-1}\mathcal{F}, f^{-1}+, f^{-1}0, f^{-1}i)$. The axioms
are going to be preserved because $f^{-1}$ is a functor
which commutes with finite limits. Finally it is not hard
to check that $f_*$ and $f^{-1}$ are adjoint as usual.

\medskip\noindent
In \cite{SGA4} this method is used. They introduce something
called an ``{\it esp\`ece the structure alg\'ebrique $\ll$d\'efinie
par limites projectives finie$\gg$}''. For such an esp\`ece you
can use the method described above to define a pair of adjoint
functors $f^{-1}$ and $f_*$ as above. This clearly works for most
algebraic structures that one encounters in practice.
Instead of formalizing this construction we simply list those
algebraic structures for which this method works (to be
verified case by case). In fact, this method works for any
morphism of topoi.

\begin{proposition}
\label{proposition-functoriality-algebraic-structures-topoi}
\begin{slogan}
Morphisms of topoi preserve algebraic structure.
\end{slogan}
Let $\mathcal{C}$, $\mathcal{D}$ be sites.
Let $f = (f^{-1}, f_*)$ be a morphism of topoi
from $\Sh(\mathcal{D}) \to \Sh(\mathcal{C})$.
The method introduced above gives rise to an adjoint
pair of functors $(f^{-1}, f_*)$ on sheaves of algebraic structures
compatible with taking the underlying sheaves of sets
for the following types of algebraic structures:
\begin{enumerate}
\item pointed sets,
\item abelian groups,
\item groups,
\item monoids,
\item rings,
\item modules over a fixed ring, and
\item lie algebras over a fixed field.
\end{enumerate}
Moreover, in each of these cases the results above labeled ($\alpha$),
($\beta$), ($\gamma$), ($\delta$), ($\epsilon$), and ($\zeta$) hold.
\end{proposition}

\begin{proof}
The final statement of the proposition holds simply since each of the listed
categories, endowed with the obvious forgetful functor, is indeed a type of
algebraic structure in the sense explained at the beginning of this section.
See Sheaves, Lemma \ref{sheaves-lemma-list-algebraic-structures}.

\medskip\noindent
Proof of (2). We think of a sheaf of abelian groups as
a quadruple $(\mathcal{F}, +, 0, i)$ as explained in the discussion preceding
the proposition.
If $(\mathcal{F}, +, 0, i)$ lives on $\mathcal{C}$, then its pullback
is defined as $(f^{-1}\mathcal{F}, f^{-1}+, f^{-1}0, f^{-1}i)$.
If $(\mathcal{G}, +, 0, i)$ lives on $\mathcal{D}$, then its pushforward
is defined as $(f_*\mathcal{G}, f_*+, f_*0, f_*i)$. This works because
$f_*(\mathcal{G} \times \mathcal{G}) = f_*\mathcal{G} \times f_*\mathcal{G}$.
Adjointness follows from adjointness of the set based functors,
since
$$
\Mor_{\textit{Ab}(\mathcal{C})}
((\mathcal{F}_1, +, 0, i), (\mathcal{F}_2, +, 0, i))
=
\left\{
\begin{matrix}
\varphi \in \Mor_{\Sh(\mathcal{C})}
(\mathcal{F}_1, \mathcal{F}_2) \\
\varphi \text{ is compatible with }+, 0, i
\end{matrix}
\right\}
$$
Details left to the reader.

\medskip\noindent
This method also works for sheaves of rings by thinking
of a sheaf of rings (with unit) as a sextuple
$(\mathcal{O}, + , 0, i, \cdot, 1)$ satisfying a list
of axioms that you can find in any elementary
algebra book.

\medskip\noindent
A sheaf of pointed sets is a pair $(\mathcal{F}, p)$, where
$\mathcal{F}$ is a sheaf of sets, and $p : * \to \mathcal{F}$
is a map of sheaves of sets.

\medskip\noindent
A sheaf of groups is given by a quadruple $(\mathcal{F}, \cdot, 1, i)$
with suitable axioms.

\medskip\noindent
A sheaf of monoids is given by a pair $(\mathcal{F}, \cdot)$
with suitable axiom.

\medskip\noindent
Let $R$ be a ring. An sheaf of $R$-modules is given by
a quintuple $(\mathcal{F}, +, 0, i, \{\lambda_r\}_{r \in R})$,
where the quadruple $(\mathcal{F}, +, 0, i)$ is a sheaf of
abelian groups as above, and $\lambda_r : \mathcal{F} \to \mathcal{F}$
is a family of morphisms of sheaves of sets
such that
$\lambda_r \circ 0 = 0$,
$\lambda_r \circ + = + \circ (\lambda_r, \lambda_r)$,
$\lambda_{r + r'} =
+ \circ \lambda_r \times \lambda_{r'} \circ (\text{id}, \text{id})$,
$\lambda_{rr'} = \lambda_r \circ \lambda_{r'}$,
$\lambda_1 = \text{id}$, $\lambda_0 = 0 \circ (\mathcal{F} \to *)$.
\end{proof}

\noindent
We will discuss the category of sheaves of modules over a sheaf of rings in
Modules on Sites, Section \ref{sites-modules-section-sheaves-modules}.

\begin{remark}
\label{remark-no-pullback-presheaves}
Let $\mathcal{C}$, $\mathcal{D}$ be sites.
Let $u : \mathcal{D} \to \mathcal{C}$ be a continuous functor
which gives rise to a morphism of sites $\mathcal{C} \to \mathcal{D}$.
Note that even in the case of abelian groups we have not defined
a pullback functor for presheaves of abelian groups.
Since all colimits are representable in
the category of abelian groups, we certainly may define
a functor $u_p^{ab}$ on abelian presheaves by the same colimits
as we have used to define $u_p$ on presheaves of sets.
It will also be the case that $u_p^{ab}$ is adjoint to
$u^p$ on the categories of abelian presheaves.
However, it will not always be the case that $u_p^{ab}$
agrees with $u_p$ on the underlying presheaves of sets.
\end{remark}












